\subsection{Peculiarities of Distributed Electric STOL Aircraft}
Objective: What about the aerodynamics and dynamic system of a STOL aircraft makes it more complicated than a normal plane
\subsubsection{Aerodynamics} \label{sec:aero}
- what is STALL? does it exist?
- control reversal / steep performance curves
- extremely coupled with thrust, lift, pitching moment
- extreme circulation
- this has ONLY to do with the 2D 3D aerodynamics, not anything about control, dynamic pressure, or any of those things

The aerodynamics and performance of blowing-augmented airfoils is much different than those for clean wings. 
The starting point for this is the understanding that while the performance of a 'clean' 2-dimensional wing with a single flap, which will define as $c_\ell, c_m, c_d$ are $f(\alpha, \delta_f, Re)$ while for a blown wing we now must consideer a new blowing parameter giving the performance as $f(\alpha,\delta_f, Tc',Re)$ \cite{spence_lift_1958}. 
Thus aerodynamics and thrust have become coupled and as \cite{courtin_performance_nodate} observes, this can make it difficult to find an operating point with net-0 x-force for landing. \\
\\
The airfoil dynamics themselves are also complex. 
Historically we prefer to fly aircraft in the linear regions of the performance curves where performance is predictable within an acceptable margin \cite{drela_flight_2014,etkin_dynamics_1959}, but for blowing-augmented aircraft this likely is not possible if we want to use them to their fullest potential. 
The performacne of these airfoils is very non-linear especially in high-lift regions as seen in \cref{fig:blown-airfoil-performance} and it seems unlikely that traditional flight controllers are suited to handling these problems.
sign reversals and large slopes can be seen in $c_m$, and even the definition of stall is hard to define exactly for these airfoils \cite{long_parametric_2021}.\\
\\

\begin{figure}{!h}
	\centering
	\includegraphics[width=0.8\textwidth]{blowing-performance.png}
	\caption{Performance curves showing the primary $c_\ell, c_x, c_m$ relationships for a blown wing as a function of $\alpha$, $\delta_f$, and $\Delta c_J$ from \cite{long_parametric_2021}. Aerodynamic performance is shown  to be extremely nonlinear along different control dimensions} \label{fig:blown-airfoil-performance}
\end{figure}

\subsubsection{Low Flight Speeds and Disturbance Sensitivity}
- one is the fact that we fly really slow
- look at courtin thesis for exact sensitivites, he did a section on this
- know from experience roll and pitch are super sensitive
- this is the motivation to look at gusts

Slow flight speeds and the nonlinear dynamics make these aircraft much more vulnerable to disturbances. While a typical Cessna 172 might land at around $60$kts tof airspeed, a blowing-augmented on might be able to do so at nearly half that $30kts$ \cite{courtin_performance_nodate}. 
Lower flight speeds make the aircraft much more sensitive to gusts in the nonlinear manner as shown in \cref{eq:gust_effect} for an aircraft in level flight encoutnering an updraft.
Thermal drafts and shears are very common near airports and more common in cities (ref needed), additionally in areas with high buildings or peaks and inconsistent ground heating much more unpredictable and much stronger than those in more consistent conditions. 
On approach the potential for large swings in $\alpha$ due to gusts, shears, and other disturbances combined with the nonlinear aerodynamics from the previous section indicate that these aircraft will be difficult to control in 'real' environments. 
\begin{equation} 
	\Delta\alpha_\mathrm{gust} = \arcsin\left(\frac{w_g}{V_\infty}\right) \label{eq:gust_effect}
\end{equation}
\begin{figure}[!h]
	\centering
	\includegraphics[width=0.8\textwidth]{generic-pic.jpeg}
	\caption{Change in $\alpha$ as a function of true airspeed. Operating points for SSTOL aircraft feel much larger effect for the same-sized gust} \label{fig:gust-effect}
\end{figure}

\subsubsection{Control Saturation} \label{sec: saturation}

%	- flying at low speeds is great, but there is tension between flying slow and having enough control authority to do anything
%	- at low speeds any unblown surface is going to have less authority
%	- additionally, with strong pitching moment, it is possible
For these aircraft designs there is a tension between the desire to blow more and fly more slowly and aircraft controllability.
The issue is that as more blowing is added to the system and the aircraft is able to achieve higher lift coefficients and thus fly more slowly, unblown surfaces on the aircraft experience both less dynamic pressure and significant influence from the circulation generated by the blown wing. 
Thus even with reasonably-sized control surfaces and throws on controls such as the elevator, it is extremely likely that errant disturbances or control inputs could cause a problematic loss of control. \\
\\
The most apparent reaosn for control saturation is the loss of dynamic pressure on unblown control surfaces. 
At high blowing setttings the velocity ratio $\frac{V_j}{V_\infty}$ have been observed to easily be as high as 3 \cite{long_parametric_2021} and will most likely be even higher for production aircraft \cite{courtin_performance_nodate}. 
It is readily apparent from the definition of dynamic pressure that the ratio between blown and unblown surfaces will obey a square law, so that the observed ratio $\frac{V_j}{V_\infty} = 3$ will lead to a dynamic pressure ratio of $\frac{q_\mathrm{blown}}{q_\mathrm{unblown}} = 9$. 
This leads to two connected issues: for short landing we want to push these aircraft to fly as slow as possible, but at some point there crossees a line where unblown controls simply do not have enough dynamic pressure and are likely to either simply not generate enough force to balance those from the wing or stall while doing so leading to the loss of control. \\
\\
This perhaps would not be a significant issue if surfaces were not signficiantly influenced by the wing, but for blown aircraft the exact opposite is true. 
While it is well known that even for conventional aircraft the wing imparts an $\Delta\alpha$ at the tail this is for the most part a very miinor impact for lift coefficients we tend to fly at (in the neighborhood of 1-2). 
Blown lift aircraft fly at much higher lift coefficients and thus from theory alone we would expect similarly high $\Delta\alpha$ around the wing from induced velocities: multiple sources confirm this to be the case \cite{courtin_performance_nodate} \cite{courtin_flight_2020} (need a few more sources here).
Large $\alpha$ offsets for blowing operating points imply that a significant control effort must be applied to simply stabilize the aircraft, most acutely along the pitch axis; such large control efforts may push these surfaces closer to stall or saturation. 
Thus it is possible to enter a situation where a small disturbance in the biased control direction could push a surface past its authority limits by either stalling it or hitting a throw limit leading to loss of control. \\
\\
Lastly, another peculiarity of blown lift aircraft is the large slopes observed in pitching moment and the sign reversal on coefficients at certain operating points. 
This is particularly apparent in the $c_m-\alpha$ curves for a 2-dimensional wing shown in \cref{fig:blown-airfoil-performance}, where for all blowing cases a steep slope and sign reversal is observed near $\alpha = 15^\circ$ for all flap deflections \cite{long_parametric_2021}. 
These control reversal can be extremely hazardous as it is possible that small perturbations can push the aircraft operating points suddenly and without warning. 
Limitations on actuation speed, as well as the bias and saturation limits described earlier in this section, mean that if these disturbances were to happen quickly enough or are large enough, it is possibel to enter an uncontrollable operating point. 
While we might seek to simply stay away from operating regions with large slopes and sign reversal, these regions also happen to overlap with high-$c_\ell$ operating regions which are desireable for STOL performance.\\
\\

\subsubsection{High-Dimensional Control Space} \label{sec:acuation}
- what actuator do you use and do respond to what?
- there are many more combiantions than previously
- control is incredibly coupled
- hard limits on some, do you use others to compensate?
For the direct control of blown-lift aircraft we can make two observations based on \cref{fig:courtin-control-layout}. The first is that the number of possible control inputs is much higher than on a conventional aircraft due to the number of propellers and flaps which can be actively managed; and the second is that each of these has a much more complex dynamic response than conventional counterparts. 
For example, a change in flap setting results in a constant shift in $C_L, C_M, C_D$ and thurst can be modulated to maintain speed. 
But based on the discussion in \cref{sec:aero}, it is quickly seen that thrust modulation will now also impact $C_L, C_D, C_M$.
The influence of this change is also extremely location dependent; flaps and propellers on outer edges of the wing have more roll impact than those more center, and the effect of induced velocity is also dramatic \cite{courtin_performance_nodate}.\\
\\
There is also the question of how-best to use the large number of contorls. 
A conventional aircraft has flight controls that each have a very clear primary axis of effect: ailerons control roll, rudder controls yaw, elevator controls pitch, throttle controls velocity, and so on. 
Many acutators gives us many possibilities and while the system is still technically under-actuated based on the definition of Chapter 1 of \cite{tedrake-underactuated}, it is very non-unique.
This suggests the question of using controls and control combinations as effectively as possible. \\

Courtin suggests control groupings as a means of handling redundency in the control space \cite{courtin_performance_nodate}.
However, it is very likely that for many maneuvers there are more effective or safe combinations of control inputs that result in the same dynamic response and with modern fly-by-wire systems control gains are easy to change.
How best explore the space of possible control inputs for different maneuvers is a question that this thesis will explore.\\
%	\cite{under-actuated robotics here}
%	over-actuated because we have many of the same control
%	- normal aircraft is nicely split so we typically have 1, maybe 2 controls per-axis for redundency
%	- think 10 propulsors so we are over-actuated on thrust
%	- but we can also see that the same applies for roll
%	-courtin tried to solve this by reducing control dimensionality and grouping controls, but this leaves out opportunity to do better
%	- non-uniqueness is the central theme, and we want to sift through the trajectory design space for 'best' solutions
\begin{figure}[!h]
	\centering
	\includegraphics[width=0.5\textwidth]{electra-el-2.png}
	\caption{Control layout for a distributed electric propulsion aircraft from Courtin} \label{fig:courtin-control-layout}
\end{figure}